% =====================================================================
%  Manuscript draft -- Synthetic SMB generation
%  Body sections for the Copernicus (gmd) manuscript.
%  Paste between \begin{document}...\end{document} of the copernicus
%  template; requires references.bib. Figure/table numbers are auto
%  via \ref{}; placeholders marked [[...]] need your input.
%
%  Verified citations are used throughout. Keys flagged VERIFY in
%  references.bib (aloni2025, mouginot2017, vandalum2024, verjans2023,
%  ultee2024, markovsky_inprep) must be confirmed before submission.
% =====================================================================

\title{Distribution-preserving spectral synthesis of surface mass balance
forcing with controllable frequency structure}

% \author[...]{...}  % fill per copernicus template

\begin{abstract}
Ice-sheet response to climate is driven not only by trends in the mean
state but by climatic \emph{noise}: variability in surface mass balance
(SMB) that spans annual to centennial timescales. Attributing the
dynamical response of an ice sheet to a specific band of that variability
is impossible from observations alone, because a single short record
confounds all timescales in one realisation. We present a method that
generates synthetic SMB series of arbitrary length that are statistically
consistent with an observed record in both their marginal distribution and
their power spectral density (PSD), and that permits controlled
amplification or suppression of variance within named frequency bands. The
method extends Fourier phase randomisation \citep{theiler1992,schreiber2000}
with three components: a Gaussian rank transform that decouples spectral
structure from the (right-skewed) marginal distribution; a semi-parametric
inverse transform that yields physically continuous extremes under
amplified forcing; and selective per-band PSD scaling that turns a
statistically consistent generator into an instrument for controlled
experiments. Calibrated on a 45-year, monthly RACMO2.4p1 record for the
Pine Island Glacier basin \citep{vandalum2025}, the method preserves the
observed mean to a ratio of $1.0000$, the variance to within $0.4\,\%$, and
the seasonal cycle to an RMSE below $3\times10^{-4}$~m~w.e.~a$^{-1}$; the
ensemble-mean PSD lies within the observed $95\,\%$ confidence interval at
every resolved frequency. An out-of-sample calibration-split test confirms
that the generator reproduces held-out statistics. Applied to PIG, the method
also shows that the basin-integrated SMB residual is statistically
indistinguishable from white noise across the annual-to-decadal range, so the
band experiments are prescribed sensitivity tests rather than perturbations of
an observed low-frequency spectrum. The same engine is
generalised to the full basin field: the spatial SMB field is decomposed
into empirical orthogonal functions (EOFs), each principal-component series
is resynthesised by the identical 1-D engine, and the synthetic modes are
recombined onto the fixed spatial patterns, yielding spatially coherent
synthetic fields whose mean and variance maps and per-mode spectra match the
observations. The method provides the forcing engine for controlled
ice-sheet experiments, in both spatially uniform and spatially resolved
form, and generalises to any geophysical field that is stationary after
simple detrending.
\end{abstract}


% #####################################################################
\introduction
% #####################################################################

The stability of marine ice sheets is governed by slow dynamics that
integrate a noisy climatic forcing. Following the classical argument of
\citet{hasselmann1976}, a slow system driven by short-correlation
forcing accumulates variance preferentially at low frequencies, so the
timescale structure of the forcing---not only its mean---shapes the
long-term response. For the Antarctic Ice Sheet, surface mass balance
(SMB) is a first-order control on this forcing, and the Pine Island
Glacier (PIG) basin is its single largest dynamic contributor to sea-level
rise \citep{rignot2019}. Understanding how the \emph{frequency structure} of
SMB variability propagates into ice dynamics is therefore central to
bounding sea-level projections.

Observed SMB records are too short to answer this question directly. A
45-year record contains only a handful of decadal cycles, so annual and
decadal variability are confounded within a single realisation, and the
record is far shorter than the multi-century response time of the ice
sheet. Isolating the dynamical contribution of a given band of SMB
variability requires forcing series that are (i) long enough for the ice
model to reach a quasi-steady state, (ii) statistically consistent with the
observations, (iii) available as independent realisations, and---critically
for controlled experiments---(iv) perturbable one frequency band at a time
while all other statistics are held fixed.

Two families of methods bear on this problem. Fourier phase randomisation
generates surrogate series that preserve the linear (second-order)
properties of a signal while randomising everything else
\citep{theiler1992,schreiber2000}; it is the natural engine for
spectrally-consistent synthesis, but in its standard form it assumes a
Gaussian marginal distribution, which SMB---right-skewed and bounded below
by zero---violates. Separately, stochastic forcing generators have been
developed specifically for ice-sheet models: \citet{muruganandham2023}
decompose a spatial ocean-forcing field into empirical orthogonal functions
(EOFs) and resynthesise the principal-component (PC) time series to produce
independent realisations of internal variability, and catchment-scale
stochastic SMB and ocean-forcing parameterisations have been built on
autoregressive models \citep[e.g.][]{ultee2024,verjans2023}. These methods
target \emph{uncertainty quantification}---sampling plausible forcing
trajectories---rather than the \emph{mechanistic isolation} of individual
timescales. \citet{muruganandham2023} in particular note the controlled
perturbation of variability bands as a direction left for future work.

This paper develops a phase-randomisation method for SMB that closes both
gaps. Our contributions are:
\begin{description}
  \item[A1 -- distribution-preserving synthesis.] A Gaussian rank transform
  decouples the marginal distribution from the spectral structure, so phase
  randomisation operates in valid Gaussian space while the synthetic series
  reproduce the observed distribution exactly.
  \item[A2 -- semi-parametric tail extrapolation.] A spliced
  empirical/parametric inverse transform allows synthetic extremes to
  exceed the observed range continuously, removing the variance-saturation
  artefact that purely empirical inverse transforms suffer under amplified
  forcing.
  \item[A3 -- frequency-band variance control.] Selective scaling of the
  PSD within named period bands amplifies or suppresses variance in those
  bands while holding the mean, the full marginal distribution, and the
  seasonal cycle fixed---converting the generator into an instrument for
  controlled experiments. This is the central advance; A1 and A2 are the
  machinery that make it valid.
  \item[A4 -- spatially coherent generalisation.] The same engine is applied
  to the full basin field through an EOF decomposition: each principal
  component is resynthesised by the identical 1-D engine and recombined onto
  fixed spatial patterns, so spatial coherence is preserved by construction
  and the basin-integrated scalar is recovered as the single-mode limit.
\end{description}
The remainder of the paper describes the data (Sect.~[[data]]), the method
component by component---including its spatial generalisation
(Sect.~[[eof]])---followed by in- and out-of-sample validation of both the
scalar and the field pipelines (Sect.~[[validation]]), a demonstration of
band control (Sect.~[[demo]]), and scope and limitations
(Sects.~[[discussion]]--[[limitations]]). An overview of the workflow, with
the 1-D core and its EOF wrapper, is given in Fig.~\ref{fig:pipeline}.


% #####################################################################
\section{Data}\label{sec:data}
% #####################################################################

We calibrate the method on basin-integrated SMB from RACMO2.4p1
\citep{vandalum2025}, the regional climate model's most recent Antarctic
product, run at 11~km horizontal resolution and forced at its lateral
boundaries by ERA5 \citep{hersbach2020}. RACMO2.4p1 introduces revised
snow-hydrometeor advection and blowing-snow physics relative to earlier
versions \citep{vandalum2024}, and its modelled SMB, energy budget and
near-surface climate compare well with observations \citep{vandalum2025}.
We use the model at its native 11~km resolution on its rotated
latitude--longitude grid. The Pine Island Glacier drainage basin is
delineated using the MEaSUREs Antarctic Boundaries product
\citep{mouginot2017}: the basin polygon is rasterised onto the RACMO grid so
that every grid cell whose centre falls inside the polygon is retained. For
the scalar (1-D) pipeline we integrate the monthly SMB over the basin
(\texttt{smbgl}) to a single series over January~1979--December~2023, giving
$n=540$ monthly values in m~w.e.~a$^{-1}$. For the field (2-D) pipeline we
retain the full masked field on the same grid, cropped to the basin bounding
box (with a two-cell margin) to reduce the domain from the full model grid
to $\approx[[n_y\times n_x]]$ cells, of which $[[n_{\mathrm{valid}}]]$ lie
inside the basin.

PIG is chosen because it is both the most demanding and the most
policy-relevant single-basin test case: it dominates present-day Antarctic
mass loss and is strongly dynamically sensitive \citep{rignot2019}. The
observed basin-mean SMB is right-skewed with a lower bound near zero
(mean~$\approx0.036$, standard deviation~$\approx0.046$~m~w.e.~a$^{-1}$;
Table~\ref{tab:data}), which directly motivates the distribution-preserving
transform of Sect.~[[gt]]. Summary statistics and a decomposition of the
record into trend, seasonal cycle and stochastic residual are shown in
Fig.~\ref{fig:data}.

Both pipelines share the same fitting and generation logic; the field
pipeline (Sect.~[[eof]]) reduces to the scalar pipeline when a single spatial
mode is retained.


% #####################################################################
\section{Methods}\label{sec:methods}
% #####################################################################

The method is organised as a fit pipeline and a generate pipeline that are
mirror images: every forward operation applied to the observations has an
inverse applied to the synthetic series (Fig.~\ref{fig:pipeline}). The
fundamental decomposition is
\begin{equation}\label{eq:decomp}
  x(t) \;=\; \mu \;+\; \tau(t) \;+\; s\!\left(m(t)\right) \;+\; r(t),
\end{equation}
where $x(t)$ is the observed monthly SMB, $\mu$ its record mean, $\tau(t)$ a
linear trend, $s(m)$ the climatological anomaly of calendar month
$m\in\{1,\dots,12\}$, and $r(t)$ the stochastic residual that is the target
of synthesis. The residual is assumed weakly stationary---constant mean,
variance and autocovariance---after removal of the deterministic terms; this
assumption is tested in Sect.~\ref{sec:pre} and probed out-of-sample in
Sect.~\ref{sec:validation}. Synthesis of $r(t)$ is performed entirely in
Gaussian space, so that the spectral structure and the marginal distribution
are manipulated independently. Two asymmetries in the pipeline are deliberate
and act as correctness guarantees. First, $\mu$ is removed once and re-added
exactly once, in the generator, so no component can drop or double-count it.
Second, $\tau(t)$ is removed at fit time but is not re-added to synthetic
series, because extrapolating a 45-year linear trend across a millennial
synthetic record is physically unjustified and would mask the variability
signal under study \citep[cf.][]{hasselmann1976}; the synthetic forcing is
therefore stationary in its mean by design.

\paragraph{Key equations.} The essential relations, referenced throughout,
are: the signal decomposition, Eq.~(\ref{eq:decomp}); the Gaussian rank
transform and its semi-parametric inverse, Eqs.~(\ref{eq:rank}) and
(\ref{eq:qsemi}); the spectral confidence interval, Eq.~(\ref{eq:ci}); the
band-scaling rule and Parseval amplitude construction, Eqs.~(\ref{eq:band})
and (\ref{eq:amp}); the phase-randomised synthesis step, Eq.~(\ref{eq:phase});
and, for the spatial pipeline, the area-weighted EOF decomposition and field
reconstruction, Eqs.~(\ref{eq:svd}) and (\ref{eq:recon2d}). Equations
(\ref{eq:qsemi})--(\ref{eq:amp}) and (\ref{eq:svd})--(\ref{eq:recon2d}) are
the load-bearing results of the paper: they define, respectively, how the
marginal distribution is preserved and controllably extrapolated, how
variance is injected into a named frequency band, and how a spatially
coherent field is built from independently-synthesised modes.

\subsection{Preprocessing: trend and seasonal removal}\label{sec:pre}

Both the Gaussian transform and the spectral estimator assume a stationary,
zero-mean input, which raw SMB is not: it carries a non-zero mean, a
deterministic annual cycle, and a possible long-term trend
\citep{vonstorch1999,wilks2011}. The preprocessor estimates the deterministic
terms of Eq.~(\ref{eq:decomp}) in order. The trend $\tau(t)$ is fitted as a
degree-1 polynomial by ordinary least squares along the time axis, and its
coefficients (not its evaluated values) are stored, so the trend
can be re-evaluated on any time coordinate. A degree-1 fit is the most
conservative defensible choice for a 45-year record: higher-order
polynomials would absorb multi-decadal variability that belongs to the
stochastic signal and would confound the decadal-band experiment of
Sect.~[[demo]]. The seasonal term $s(m)$ is then estimated as the mean for
each of the twelve calendar months, computed \emph{on the detrended series}
so that a residual trend does not bias the climatology, and subtracted. An
empirical monthly climatology is preferred over harmonic regression because
it makes no shape assumption and each month is well-sampled ($45$
observations). Removing the seasonal cycle reduces the power at the annual
frequency by a factor of $\approx30$, leaving in the residual only the
stochastic year-to-year variation in seasonal amplitude, which is the
physically meaningful quantity to synthesise. Stationarity of the residuals
is confirmed with an augmented Dickey--Fuller test
\citep{dickey1979,said1984}, which rejects the unit-root null at the
$5\,\%$ level; we note that the ADF test has limited power at $n=540$ and
provides a necessary but not sufficient condition for stationarity.

\subsection{Gaussian rank transform (A1)}\label{sec:gt}

Phase randomisation requires a Gaussian input. We obtain one by a rank
transform: the residuals $r_1,\dots,r_n$ are sorted, the $i$-th order
statistic $r_{(i)}$ is assigned the Hazen plotting position
\citep{hazen1914}, and each value is mapped to the corresponding
standard-normal quantile,
\begin{equation}\label{eq:rank}
  p_i=\frac{i-\tfrac{1}{2}}{n},
  \qquad
  g_{(i)}=\Phi^{-1}(p_i),
  \qquad i=1,\dots,n,
\end{equation}
where $\Phi$ is the standard-normal CDF. The transformed series takes
exactly the values $\{\Phi^{-1}(p_i)\}$, i.e.\ its empirical distribution
matches the standard normal at every quantile by construction (a one-sample
Kolmogorov--Smirnov test on the PIG record gives $D=0.0009$); the Hazen
positions are strictly interior to $(0,1)$, so all quantiles are finite.
Because the map from ranks to values is monotone, the transform preserves
the ordering, and therefore the rank-correlation structure, of the
residuals. This is the probability-integral
transform in rank form \citep{rosenblatt1952}, applied here specifically to
decouple the marginal distribution from the spectrum so that phase
randomisation in Gaussian space does not distort the distribution
\citep{theiler1992}. No subsequent step alters the marginal distribution
except the inverse of this transform.

\subsection{Semi-parametric inverse transform (A2)}\label{sec:inv}

Back-transformation maps a synthetic Gaussian series to the physical
distribution: $u=\Phi(g)$ (clipped to $[10^{-15},1-10^{-15}]$), followed by
a semi-parametric inverse cumulative distribution function (CDF),
\begin{equation}\label{eq:qsemi}
  Q_{\mathrm{semi}}(u)=
  \begin{cases}
    Q_{\mathcal{N}}(u;\hat\mu_t,\hat\sigma_t)+\delta_{\mathrm{low}},
      & u < p_1,\\[2pt]
    \mathrm{interp}\!\left(u;\,\{p_i\},\{r_{(i)}\}\right),
      & p_1 \le u \le p_n,\\[2pt]
    Q_{\mathcal{N}}(u;\hat\mu_t,\hat\sigma_t)+\delta_{\mathrm{high}},
      & u > p_n,
  \end{cases}
\end{equation}
where $Q_{\mathcal{N}}$ is the quantile function of a normal distribution
fitted to the residuals by maximum likelihood, and the continuity offsets
\begin{equation}\label{eq:delta}
  \delta_{\mathrm{low}} = r_{(1)}-Q_{\mathcal{N}}(p_1),
  \qquad
  \delta_{\mathrm{high}} = r_{(n)}-Q_{\mathcal{N}}(p_n)
\end{equation}
make the spliced function continuous at both boundaries by construction
(verified numerically to $<10^{-10}$). Within the observed range the
synthetic marginal therefore follows the empirical distribution of the
training residuals, interpolating linearly between order statistics; beyond
it, extremes follow the shifted normal tail.

The extrapolation is essential for the band experiments. A purely empirical
inverse CDF clips every value to the observed minimum and maximum, so under
$5\times$--$10\times$ band amplification the extremes pile up at those
bounds and the synthetic variance \emph{saturates} rather than scaling with
the imposed factor---conflating spectral structure with an artificial
variance ceiling and rendering the controlled experiments
uninterpretable. The parametric tail removes this ceiling; the splice
points sit at the exact empirical boundaries, so the round-trip is exact
for all training data (maximum pointwise error $<3\times10^{-9}$). A normal
tail is a deliberate simplification: for right-skewed residuals it
underestimates large positive anomalies, and a generalised Pareto tail
fitted to exceedances would be the principled alternative
\citep{coles2001}; we adopt the normal tail because the band experiments
already extrapolate into a regime where the exact tail shape is a secondary
concern, and we flag this in Sect.~[[limitations]]. Finally, because passing
a symmetric Gaussian through the asymmetric map of Eq.~(\ref{eq:qsemi})
produces a non-zero sample mean---large enough under $10\times$ scaling to
masquerade as a change in mean accumulation---the residual mean is
subtracted inside the inverse transform, guaranteeing that every synthetic
series has exactly zero mean before $\mu$ is restored.

\subsection{Spectral estimation}\label{sec:psd}

The one-sided PSD of the Gaussianised residuals is estimated by Welch's
method \citep{welch1967}: segments of $\texttt{nperseg}=60$ months (5~yr)
with $50\,\%$ overlap and a Hann window, giving $K=17$ averaged
periodograms at $31$ frequencies. Segment averaging trades frequency
resolution for variance reduction. The segment length sets the Rayleigh
frequency $\Delta f = f_s/\texttt{nperseg} = 0.2$~cycles~yr$^{-1}$, so the
longest independently resolved period is 5~yr: the annual band
($[0.8,1.5]$~yr, i.e.\ $[0.67,1.25]$~cycles~yr$^{-1}$) is covered by three
resolved bins, whereas the decadal band ($[8,20]$~yr, i.e.\
$[0.05,0.125]$~cycles~yr$^{-1}$) lies entirely \emph{below} $\Delta f$ and
contains no independently resolved bin. Power assigned to the decadal band
during synthesis is therefore obtained by interpolation between the
zero-frequency term and the lowest resolved bin, rather than estimated from
data at those frequencies; the consequences for the interpretation of the
decadal experiment are set out in Sect.~[[limitations]]. The $(1-\alpha)$
confidence interval
on the estimate $\hat P(f)$ follows the chi-squared sampling distribution of
an averaged periodogram \citep{percival1993},
\begin{equation}\label{eq:ci}
  \left[\;
    \frac{\nu\,\hat P(f)}{\chi^2_{\nu,\,1-\alpha/2}},\;\;
    \frac{\nu\,\hat P(f)}{\chi^2_{\nu,\,\alpha/2}}
  \;\right],
  \qquad \nu = 2K = 34,
\end{equation}
which treats the $K$ windowed segments as independent; the $50\,\%$ overlap
correlates neighbouring segments, so Eq.~(\ref{eq:ci}) is an approximation,
and the resulting $\approx2.6\times$ ratio between the
upper and lower bounds is the honest statement of spectral uncertainty
imposed by a 45-year record. The observed PSD and its $95\,\%$ interval are
shown in Fig.~\ref{fig:psd}.

The choice $\texttt{nperseg}=60$ is supported empirically. Because Welch's
method removes the mean of each segment, variance carried by fluctuations
\emph{between} segments---that is, at periods longer than the segment---is
not represented in the estimate. For the PIG residual this loss is
negligible: the $\texttt{nperseg}=60$ estimate accounts for $99.7\,\%$ of the
total residual variance (Parseval check), so essentially no power is being
discarded at low frequencies. The reason is that the residual is close to
white. Its lag-1 autocorrelation is $\phi=0.08$ (an $e$-folding memory of
$0.4$~months), and band-integrated multitaper estimates on the full record
(NW${}=1.5$, resolution bandwidth $0.067$~cycles~yr$^{-1}$) place only
$1.7\,\%$ of the residual variance in the decadal band and $11.5\,\%$ in the
annual band, against white-noise expectations of $1.25\,\%$ and $9.7\,\%$
respectively. Neither excess is distinguishable from a white-noise or an
AR(1) null (Monte Carlo, 500 surrogates; observed at the 84th and 78th
percentiles respectively). Within the resolution of a 45-year record, PIG
basin-integrated SMB variability is therefore statistically indistinguishable
from white noise across the annual-to-decadal range. This is itself a
result---and it is precisely the forcing regime for which the
noise-integration argument of \citet{hasselmann1976} predicts that a slow
system will generate red-spectrum variability from white-spectrum forcing.

\subsection{Phase randomisation, band scaling and ensemble synthesis
(A3)}\label{sec:synth}

To synthesise a series of $N$ months, the fitted PSD is
interpolated onto the synthetic frequency grid (for a 1000-year series,
$6001$ points against the $31$ fitted). Band scaling is applied here:
within a target period band $[p_{\min},p_{\max}]$ the interpolated PSD is
multiplied by a factor $\gamma_b$,
\begin{equation}\label{eq:band}
  \tilde P(f)=
  \begin{cases}
    \gamma_b\,\hat P(f), & 1/p_{\max}\le f \le 1/p_{\min},\\[2pt]
    \hat P(f), & \text{otherwise},
  \end{cases}
\end{equation}
which, by Parseval's theorem, scales the Gaussian-space variance contributed
by that band by $\gamma_b$ and the band standard deviation by
$\sqrt{\gamma_b}$. Because the inverse transform of Eq.~(\ref{eq:qsemi}) is
monotone but nonlinear, the corresponding scaling of the \emph{physical}
band variance is approximate rather than exact; the realised scaling is
quantified empirically in Sect.~[[demo]]. Equation~(\ref{eq:band}) is the
operational definition of a controlled experiment (A3) and follows the
band-perturbation direction that \citet{muruganandham2023} identify as
future work. Fourier amplitudes are constructed from the (scaled) PSD as
\begin{equation}\label{eq:amp}
  A_k=\sqrt{\tilde P(f_k)\,\Delta f\,N},
\end{equation}
which is the discrete Parseval-consistent normalisation
\citep{percival1993}: summing $|A_k|^2$ over frequencies recovers the target
variance. Each ensemble member $j$ then draws independent phases and inverts
to the time domain,
\begin{equation}\label{eq:phase}
  Z_k^{(j)}=A_k\,e^{i\phi_k^{(j)}},\qquad
  \phi_k^{(j)}\sim U(-\pi,\pi),\qquad
  g^{(j)}=\mathrm{irfft}\!\left(Z^{(j)}\right)\big/\hat\sigma^{(j)},
\end{equation}
with the zero-frequency and Nyquist phases fixed at zero so that the inverse
transform is real-valued, and $\hat\sigma^{(j)}$ the sample standard
deviation of the inverted series, enforcing exact unit variance per member.

Two properties follow directly. Because the amplitude spectrum is
deterministic, $\mathbb{E}\,|Z_k|^2=A_k^2$ independently of the phase draw,
so the ensemble-mean PSD converges to the target
\citep{theiler1992,schreiber2000}; and because phases are drawn
independently across members, the expected cross-covariance between any two
members vanishes at all lags (sample cross-covariances fluctuate around zero
at the $O(N^{-1/2})$ level expected for independent series). The
unit-variance normalisation in Eq.~(\ref{eq:phase}) removes member-to-member
variance drift introduced by interpolation onto the fine grid without
altering the spectral shape (it is a scalar multiplication).

\subsection{Reconstruction and experiment configuration}\label{sec:recon}

The synthetic Gaussian series is passed back through the inverse transform
(Sect.~\ref{sec:inv}), the observed mean is added once, and the seasonal
cycle is restored; the trend is not restored (Sect.~\ref{sec:pre}). Each
ensemble is fully specified by an \texttt{Experiment} configuration---the
number of years and members, a random seed, and the set of band scale
factors---which is serialised alongside the output so that every synthetic
dataset is reproducible from its configuration alone (RACMO input
$\rightarrow$ configuration $\rightarrow$ ensemble). Within an experiment
suite the members share a random seed across band conditions, so that any
downstream difference is attributable to the band scaling rather than to the
phase draws; independent seeds are used when quantifying ensemble spread.

\subsection{Spatial generalisation via EOF decomposition (A4)}\label{sec:eof}

The scalar engine above generalises to the full basin field by inserting an
EOF decomposition before the spectral engine and an EOF recomposition after
it; the 1-D Gaussian transform, spectral synthesiser and experiment
configuration are reused without modification, and applied once per retained
mode.

\subsubsection{Spatial preprocessing}
Trend and seasonal cycle are removed \emph{per grid cell}. A degree-1
polynomial is fitted independently at each cell (over an integer time index,
to avoid calendar-arithmetic artefacts), the per-cell time mean is stored
separately, and a per-cell monthly climatology is computed on the detrended,
mean-centred field. Removing the seasonal cycle in physical space, before the
decomposition, is deliberate: the seasonal cycle of SMB is itself spatially
structured (its amplitude and phase vary with orography and distance from the
coast), so leaving it in would cause the leading EOF to capture the spatial
pattern of the \emph{seasonal cycle} rather than the stochastic variability
of interest. As in the scalar case, the per-cell mean is restored during
reconstruction but the per-cell trend is not.

\subsubsection{EOF decomposition}
The residual field is flattened to a matrix
$\mathbf{X}\in\mathbb{R}^{T\times P}$ over the $P$ valid (in-basin) cells and
area-weighted per cell $j$ by $w_j=\sqrt{\cos\phi_j}$ ($\phi_j$ the cell
latitude), so that the decomposition minimises area-weighted rather than
cell-count variance \citep{north1982,hannachi2007}. An economy singular value
decomposition of the weighted matrix then defines the modes,
\begin{equation}\label{eq:svd}
  \mathbf{X}_w=\mathbf{X}\,\mathrm{diag}(w)
  =\mathbf{U}\mathbf{S}\mathbf{V}^{\top},
  \qquad
  \mathbf{E}=\mathbf{V}^{\top},\quad
  \mathbf{P}=\mathbf{U}\mathbf{S},\quad
  \lambda_k=\frac{s_k^2}{\sum_j s_j^2},
\end{equation}
where the rows of $\mathbf{E}$ are the orthonormal
spatial patterns (EOFs), the columns of $\mathbf{P}$ the principal-component
series, and $\lambda_k$ the fraction of area-weighted variance explained by
mode $k$ \citep{lorenz1956,preisendorfer1988}. EOF
signs are standardised (first non-zero loading positive) for
interpretability. The decomposition is truncated at $K$ modes, chosen as the
smallest $K$ whose cumulative explained variance reaches a target threshold
(the implementation reports the $80\,\%$ and $95\,\%$ thresholds and defaults
to $K=[[K]]$ modes, explaining $[[\text{cumvar}]]\%$ of the area-weighted
residual variance for PIG; see the scree plot, Fig.~\ref{fig:scree}). North's
rule of thumb \citep{north1982} bounds the modes that are well separated from
their neighbours given the effective sample size, and is used to flag modes
whose patterns should not be over-interpreted.

\subsubsection{Per-mode synthesis and recombination}
Each retained PC is a scalar, zero-mean series and is passed through the
identical 1-D pipeline: the Gaussian rank transform (Sect.~\ref{sec:gt}), the
Welch spectral fit (Sect.~\ref{sec:psd}), and phase-randomised synthesis with
optional band scaling (Sect.~\ref{sec:synth}). Modes with negligible variance
are skipped. For an ensemble member, each (member, mode) pair is synthesised
with an independent random-number stream, so that modes are mutually
independent and members are independent, while the whole ensemble remains
reproducible from the experiment seed. The synthetic PCs are back-transformed
to physical units, recombined onto the fixed spatial patterns, and
un-weighted,
\begin{equation}\label{eq:recon2d}
  \tilde r(\mathbf{x}_j,t)
  = \frac{1}{w_j}\sum_{k=1}^{K}\widetilde{PC}_k(t)\,E_k(\mathbf{x}_j),
\end{equation}
after which the per-cell seasonal cycle and mean are restored,
giving a synthetic field ensemble with dimensions
(member, time, $y$, $x$). Because the spatial patterns are held fixed and only
the temporal PCs are resynthesised, the synthetic field reproduces the
observed spatial covariance structure of the retained modes by
construction---the same principle
\citet{muruganandham2023} use for ocean forcing, here applied to atmospheric
SMB with the distribution-preserving, band-controllable 1-D engine acting on
each mode. Band scaling (A3) applied to a mode perturbs the temporal
variability of a coherent spatial pattern, not of individual grid cells. The
basin-integrated scalar of Sects.~\ref{sec:pre}--\ref{sec:recon} is the
single-mode limit of this construction.


% #####################################################################
\section{Validation}\label{sec:validation}
% #####################################################################

We distinguish in-sample fidelity---necessary but circular, since the
generator is fitted and scored on the same record---from out-of-sample
generalisation, which is the primary defence against overfitting.

\subsection{In-sample fidelity}

Across the standard experiments the generator preserves the deterministic
components to numerical precision: the mean ratio (synthetic/observed,
pooled over members and time) is
$1.0000$ for all experiments, including $10\times$ band scaling; the
variance ratio is $\approx0.996$; and the seasonal-cycle RMSE is
$\approx2\times10^{-4}$~m~w.e.~a$^{-1}$, i.e.\ floating-point accumulation
error (Table~\ref{tab:metrics}). Mean preservation is exact by construction
(Sect.~\ref{sec:inv}); the small variance deficit arises from linear
interpolation between order statistics in Eq.~(\ref{eq:qsemi}), which
slightly smooths the empirical distribution. Spectral and temporal fidelity
are measured
by the PSD RMS error and the lag-1 and lag-12 autocorrelation errors; for
band-scaled experiments the PSD error is \emph{expected} to rise in the
scaled band, and that rise is part of the intended perturbation rather than
a defect (Sect.~[[demo]]).

\subsection{Spectral coverage}

The strongest in-sample result is spectral coverage: the ensemble-mean PSD,
computed with the same Welch parameters as the fit, lies within the fitted
$95\,\%$ confidence interval at all $31$ frequencies (coverage $=1.000$),
with the synthetic/observed PSD ratio falling in $[0.68,1.33]$ across the
band (Fig.~\ref{fig:valsuite}a). The synthetic ensemble thus reproduces the
observed spectrum within its own estimated uncertainty.

\subsection{Out-of-sample calibration-split test}

To break the circularity of in-sample scoring we split the record at its
midpoint ($\approx$2001), fit an independent generator on each half, and
score each against the held-out half in both directions
\citep{michaelsen1987}. The mean ratio is $1.19$ and $0.84$ and the variance
ratio $1.29$ and $0.78$ in the two directions (Table~\ref{tab:split}). These
departures from unity are not pipeline failures: they reflect a genuine
difference in mean and variance between the two halves of the record---the
same non-stationarity that motivates the trend removal in
Sect.~\ref{sec:pre}. The test demonstrates that spectral structure fitted on
one half suffices to reproduce the statistics of the other, which is the
direct answer to the objection that the synthetic series are merely the
observed data with randomised phases.

\subsection{Ensemble convergence and representativeness}

The variance ratio is stable at $\approx0.995$ from as few as five members,
and the mean ratio is exactly $1.0$ for all ensemble sizes; we adopt $30$
members as a conservative choice for reliable tail and return-period
estimates rather than for mean or variance convergence
(Table~\ref{tab:convergence}). Finally, drawing 45-year windows from a
1000-year ensemble places the observed 45-year statistics within the
synthetic window distribution (Fig.~\ref{fig:windows}), which reframes the
question ``does the synthetic match the observations?'' as the stronger
``is the observed record a typical realisation of the fitted process?''.

\subsection{Validation of the spatial field}\label{sec:spatialval}

The field pipeline is validated with the spatial analogue of the scalar
suite, reusing the 1-D validator on each mode. Four questions are addressed.
First, does the synthetic ensemble preserve the observed \emph{spatial mean
field}? The Pearson correlation between the observed and synthetic per-cell
time-mean maps is $[[\text{mean\_map\_corr}]]$, and the global mean ratio is
$[[\text{mean\_ratio\_2D}]]$. Second, is the \emph{spatial distribution of
variance} reproduced? The correlation between the observed and synthetic
per-cell variance maps is $[[\text{var\_map\_corr}]]$ with a domain-mean
variance ratio of $[[\text{var\_map\_ratio}]]$; the observed, synthetic and
ratio maps are shown in Fig.~\ref{fig:varmaps}. Third, do the \emph{modes
carry the correct temporal spectrum}? Running the scalar spectral check on
each PC shows that the Gaussianised observed PC spectrum and the fitted PSD
agree within the $95\,\%$ confidence band for the leading modes
(Fig.~\ref{fig:pcpsd}); the $K$ retained modes capture
$[[\text{eof\_cumvar}]]\%$ of the area-weighted residual variance. Fourth,
does the field pipeline \emph{generalise out of sample}? A spatial
calibration-split test---fitting the full field pipeline (preprocessing, EOF,
per-mode synthesis) on one half of the record and scoring the synthetic field
against the held-out half in both directions---returns mean-map correlations
of $[[\text{split\_2D}]]$, with the same trend-driven mean and variance
asymmetry seen in the scalar test. Together these confirm that the synthetic
fields reproduce the observed spatial covariance structure of the retained
modes, not merely the basin-mean statistics.


% #####################################################################
\section{Demonstration of frequency-band control}\label{sec:demo}
% #####################################################################

We finally verify that band scaling produces clean perturbations at the SMB
level, independent of any ice-dynamic model. We define an annual band of
$[0.8,1.5]$~yr and a decadal band of $[8,20]$~yr. Because the observed
residual is statistically indistinguishable from white noise
(Sect.~\ref{sec:psd}), these experiments are \emph{prescribed sensitivity}
experiments: they impose a specified redistribution of variance across
timescales and ask what an ice sheet does with it. They are not perturbations
of an observationally constrained decadal spectrum, and we do not claim them
to be.

Amplifying a band multiplies the synthetic Gaussian-space power there by the
imposed factor $\gamma_b$ (Eq.~\ref{eq:band}; Fig.~\ref{fig:bands}a). The
per-member normalisation to unit variance in Eq.~(\ref{eq:phase}) means that
band scaling \emph{redistributes} variance rather than adding it. Writing
$s_b$ for the baseline fraction of variance in band $b$, the post-scaling
variance ratios are
\begin{equation}\label{eq:redistrib}
  \frac{\sigma^2_{\mathrm{in}}(\gamma_b)}{\sigma^2_{\mathrm{in}}(1)}
  = \frac{\gamma_b}{1+(\gamma_b-1)s_b},
  \qquad
  \frac{\sigma^2_{\mathrm{out}}(\gamma_b)}{\sigma^2_{\mathrm{out}}(1)}
  = \frac{1}{1+(\gamma_b-1)s_b},
\end{equation}
so total variance is conserved, in-band variance rises by less than
$\gamma_b$, and out-of-band variance necessarily falls. The invariants that
do hold exactly are the mean (Sect.~\ref{sec:inv}), the total variance, the
full marginal distribution (KS statistic unchanged) and the seasonal cycle
(Table~\ref{tab:bands}).

Equation~(\ref{eq:redistrib}) has a consequence for experiment design that is
easy to miss: the same nominal factor buys very different perturbations in
the two bands, because their baseline shares differ by nearly an order of
magnitude ($s_{\mathrm{dec}}=0.017$, $s_{\mathrm{ann}}=0.115$). At
$\gamma_b=10$ the decadal band reaches $14\,\%$ of total variance while
out-of-band variance falls by $13\,\%$; the annual band reaches $57\,\%$
while out-of-band variance falls by $51\,\%$. Making the decadal band carry
half the total variance requires $\gamma_b\approx59$, against
$\gamma_b\approx8$ for the annual band. Experiments intended to be compared
against one another should therefore be matched on the realised variance
partition rather than on the nominal factor, and we report both
(Table~\ref{tab:bands}). With that matching, experiments differ only in how
variance is distributed across timescales---the interpretive contract that
the companion study relies on \citep{markovsky_inprep}.


% #####################################################################
\section{Discussion}\label{sec:discussion}
% #####################################################################

The method supplies long, stationary, statistically-consistent and
independently-replicated SMB forcing whose spectral content can be
controlled band by band---the four requirements set out in the
introduction, plus the fifth that distinguishes it from prior work. It
extends Fourier phase randomisation \citep{theiler1992,schreiber2000} with a
distribution-preserving transform (A1) and continuous tails (A2), and it
operationalises the band-perturbation direction that \citet{muruganandham2023}
left open (A3), applied to atmospheric SMB rather than ocean forcing. Unlike
autoregressive stochastic parameterisations aimed at uncertainty
quantification \citep{ultee2024,verjans2023}, it is designed for the
mechanistic isolation of individual timescales. Nothing in the construction
is specific to PIG or to SMB: it applies to any scalar geophysical series
that is stationary after simple detrending, and, through the EOF extension
(Sect.~\ref{sec:eof}), to spatially coherent fields. The method controls the
forcing only; it makes no claim about ice response, which is the subject of
the companion paper \citep{markovsky_inprep}.

The PIG application also clarifies what the band experiments mean. Finding
the observed residual to be near-white is not a negative result: it is the
forcing regime in which \citet{hasselmann1976} predicts that a slow system
generates red-spectrum variability by integration alone, with no
corresponding structure in the forcing. Any decadal signal in the modelled
ice response to baseline forcing is therefore attributable to the ice
dynamics rather than inherited from the SMB. Against that baseline, the
band-scaled experiments become explicit counterfactuals---if SMB variability
were concentrated at annual or at decadal timescales, with the same total
variance and the same marginal distribution, how would the response differ?
That is a sharper question than perturbing an already-structured spectrum,
and it is answerable precisely because the baseline is unstructured.


% #####################################################################
\section{Limitations}\label{sec:limitations}
% #####################################################################

Several caveats bound the method's applicability.
\emph{Stationarity.} The Welch PSD and empirical CDF assume the 1979--2023
record is stationary; the calibration-split asymmetry shows it is not
strictly so, and the fitted spectral model is a temporal average. Linear
detrending addresses only first-order non-stationarity.
\emph{Decadal resolution.} At $\texttt{nperseg}=60$ the Rayleigh frequency is
$0.2$~cycles~yr$^{-1}$, so the decadal band
($[0.05,0.125]$~cycles~yr$^{-1}$) contains no independently resolved
spectral estimate; its power is set by interpolation between the
zero-frequency term and the lowest resolved bin. Two independent checks show
that little is lost by this: the estimate accounts for $99.7\,\%$ of the
residual variance, and full-record multitaper estimates find no decadal
excess distinguishable from a white or AR(1) null (Sect.~\ref{sec:psd}). The
segment length is therefore not the binding constraint---the record is. A
45-year series contains at most four or five realisations of a decadal
cycle, which is too few to constrain decadal power regardless of estimator;
$\texttt{nperseg}=240$ would resolve the band but with only $\nu\approx6$
degrees of freedom. The consequence is interpretive rather than technical:
the decadal experiment imposes variance that the observations cannot
constrain, and must be read as a prescribed sensitivity test
(Sect.~\ref{sec:demo}) rather than as a perturbation of an observed decadal
spectrum. We prefer to state this plainly rather than to imply a constraint
the data cannot supply.
\emph{Deterministic amplitude spectrum.} All members share one amplitude
spectrum and differ only in phase, so between-member variability
underestimates true amplitude variability; an amplitude-randomising
extension is a natural next step.
\emph{Tail model.} The normal tail underestimates right-skewed extremes
beyond the observed range; a generalised Pareto tail \citep{coles2001} is
the principled upgrade.
\emph{Band edges and boundaries.} Band scaling is applied as an untapered
step at the band edges, and the annual and decadal band definitions,
although physically motivated, are not unique; a tapered transition and a
boundary-sensitivity test are straightforward refinements.
\emph{EOF truncation and mode assumptions.} The field pipeline retains only
the leading $K$ modes, so variance in the discarded high-order modes---often
the small-scale, less spatially coherent component---is not reproduced;
$K$ trades spatial fidelity against the number of PC series that must be
well-fitted. The synthetic spatial patterns are exactly the observed EOFs and
are held fixed, so the method reproduces the observed modes of variability but
cannot generate spatial patterns absent from the training field. Each PC is
synthesised independently, which reproduces the marginal spectrum of every
mode but discards any cross-mode phase coupling (e.g.\ propagating features
represented jointly by two modes); and EOFs near-degenerate in variance are
not robustly separable \citep{north1982}, so individual high-order patterns
should not be over-interpreted. The per-cell mean and seasonal cycle are
assumed stationary across the synthetic horizon, as in the scalar case.


% #####################################################################
\conclusions
% #####################################################################

We have presented a method for generating synthetic SMB forcing that
reproduces the observed marginal distribution and power spectrum of a short
record and permits controlled perturbation of individual frequency bands.
The three advances---distribution-preserving synthesis (A1), semi-parametric
tail extrapolation (A2), and frequency-band variance control (A3)---together
convert Fourier phase randomisation into an instrument for controlled
ice-sheet forcing experiments, and a fourth (A4) generalises the identical
engine to spatially coherent fields through an EOF decomposition. Calibrated
on RACMO2.4p1 SMB for the Pine Island Glacier basin \citep{vandalum2025}, the
generator preserves the mean exactly and reproduces the observed spectrum
within its $95\,\%$ confidence interval, and it passes an out-of-sample
calibration-split test. Band scaling perturbs a named band while holding the
mean, marginal distribution and seasonal cycle fixed, giving experiments that
differ only in spectral structure. Applied per mode, the same engine produces
synthetic fields that reproduce the observed spatial mean and variance
patterns and per-mode spectra while preserving spatial coherence by
construction. The generator is the forcing engine for the ice-dynamic
experiments of the companion study \citep{markovsky_inprep}, in both
spatially uniform and spatially resolved form, and generalises to other
scalar and gridded climate series that are stationary after simple
detrending.


% #####################################################################
%\codeavailability{ ... Zenodo DOI ... }
%\dataavailability{ RACMO2.4p1: \citet{vandalum2025}; generated ensembles archived with Experiment metadata at [[DOI]]. }
%\appendix
%\section{Semi-parametric inverse CDF and Welch CI derivations}\label{app:eqs}
%\section{Sensitivity analyses (nperseg, band boundaries, plotting positions, tail model)}\label{app:sens}

\bibliographystyle{copernicus}
\bibliography{references}
